% -----------------------------------------------------------------------------
\section{Introduction}
\label{sec:intro}
% -----------------------------------------------------------------------------

This document defines the current test plan for \textbf{MoonAI}, a modular predator-prey
simulation platform used to study continuous neuroevolution with NEAT. The report is
written against the repository state that exists today and describes 
testing assets, workflows, and verification activities that are implemented in the
codebase, build system, or project infrastructure.

MoonAI contains several verification-relevant subsystems: C++ simulation and evolution
libraries, an SFML-based visualization frontend, an optional CUDA acceleration path,
command-line experiment orchestration, structured output logging, a Python analysis
pipeline, a profiler executable, and GitHub Actions workflows for automated build and test
execution. The test plan therefore has to cover more than unit-level algorithm checks; it
must also describe end-to-end execution, artifact generation, performance observation, and
manual user-oriented smoke testing.

\subsection{Purpose}

The purpose of this test plan is to:

\begin{itemize}[nosep]
	\item identify the MoonAI items currently under test,
	\item describe the features that are verified by automated and manual checks,
	\item define the test methodologies that are already implemented,
	\item record test control procedures, schedule, risks, and deliverables.
\end{itemize}

\subsection{Document Baseline}

The report is based on the following implemented project assets.

\begin{itemize}[leftmargin=*]
	\item \textbf{\texttt{tests/test\_evolution.cpp}.} Provides automated GoogleTest coverage
	      for NEAT data structures, mutation, crossover, speciation, JSON round-trip behavior,
	      and neural network activation.

	\item \textbf{\texttt{tests/test\_simulation\_sensors.cpp}.} Provides automated
	      GoogleTest coverage for simulation sensor encoding and density calculations.

	\item \textbf{\texttt{tests/CMakeLists.txt}.} Registers the test executable with
	      GoogleTest discovery so the implemented unit tests can be executed through CTest.

	\item \textbf{\texttt{just test} / \texttt{ctest}.} Serve as the standard local entry
	      point for automated regression execution.

	\item \textbf{\texttt{.github/workflows/ci.yml}.} Executes cross-platform automated
	      build and test runs on Linux and Windows in Debug and Release configurations.

	\item \textbf{\texttt{.github/workflows/release.yml}.} Verifies release builds by
	      running the automated test suite before packaging.

	\item \textbf{\texttt{src/main.cpp} and \texttt{src/app.cpp}.} Provide the runtime path
	      used to verify configuration validation, CLI orchestration, headless or GUI execution,
	      CPU or GPU selection, and output generation.

	\item \textbf{\texttt{analysis/moonai\_analysis/}.} Implements the post-run analysis
	      pipeline that converts experiment outputs into a self-contained HTML report.

	\item \textbf{\texttt{profiler/moonai\_profiler/} and \texttt{src/profiler\_main.cpp}.}
	      Implement the built-in performance test and timing-report path for frame-time and
	      scoped performance inspection.
\end{itemize}

\subsection{Goals}

The implemented testing approach focuses on the following goals:

\begin{itemize}[nosep]
	\item correctness of the highest-risk algorithmic logic in evolution and sensing,
	\item safe handling of configuration values and command-line overrides,
	\item stable end-to-end execution in GUI, headless, CPU-only, and GPU-enabled modes,
	\item integrity of generated output artifacts such as \texttt{config.json},
	      \texttt{stats.csv}, \texttt{species.csv}, and \texttt{genomes.json},
	\item repeatable build and regression execution on both Linux and Windows,
	\item visibility into runtime performance through the profiler and Nsight tooling,
	\item internal user-oriented acceptance checks before accepting release-candidate builds.
\end{itemize}
